\begin{theorem}[A six-spin operation moving one plaquette flux]%
    \label{thm:peps_torus_two_plaquette_physical_flux_move}
    \lean{TNLean.PEPS.IsGIsometric.exists_unitary_torusTwoPlaquettePhysicalFluxMove}
    \leanok
    \uses{thm:peps_two_plaquette_torus_geometry,
        thm:peps_torus_two_plaquette_flux_holonomy,
        thm:peps_regular_two_cycle_physical_flux_move}
    Let the horizontal and vertical torus periods be at least four and three,
    respectively. Let the original four-leg tensor be $G$-isometric for the
    regular virtual action. In the actual six-site region
    $R=\{0,1,2\}\times\{0,1\}$, let $u_g$ insert $g$ on the upward bond at
    horizontal coordinate zero, with the identity on every other bond.
    Let $u'_g$ insert $g$ on both upward bonds at horizontal coordinates zero
    and one, again leaving every other bond unchanged.
    There exists one unitary $W$ on the six original physical spins such that
    \begin{align}
        W\Psi_R(u_g,\theta)=\Psi_R(u'_g,\theta)
    \end{align}
    for every $g\in G$ and every crossing configuration $\theta$.
    Here both vectors are open-region contractions of the original tensor.
    The same unitary works for every $g$ and $\theta$.
    The left and right plaquette holonomies change from $(g^{-1},1)$ to
    $(1,g^{-1})$.
    This is the regular-representation six-site specialization of the local
    construction in \cite[Theorem~6.16]{Schuch2010PEPS}. It makes no
    parent-Hamiltonian ground-space assertion.
\end{theorem}

\begin{proof}\leanok
    Four-leg $G$-isometry gives the same local isometry in incident-edge
    coordinates. The fixed five-edge path through the six sites is a spanning
    tree; its two missing bonds are exactly the leftmost and middle upward
    vertical bonds. Apply the physical two-cycle operation to this fixed tree,
    its fixed root, and these two distinct bonds. Identifying its assignments
    with the literal bond insertions gives the asserted identity on every
    crossing column. The actual plaquette holonomy calculation gives the
    stated movement of the flux.
\end{proof}
